\chapter{Respiratory Distress}

\section*{Definition}
Respiratory distress is a clinical state of increased work of breathing, defined by the presence of tachypnea, use of accessory muscles, nasal flaring, intercostal/subcostal retractions, and possibly grunting or paradoxical breathing. It precedes respiratory failure and is graded by the degree of compensatory effort.

\section*{What Happens in Your Body}
Respiratory distress represents the body's attempt to maintain adequate ventilation and oxygenation when pulmonary mechanics are compromised. Increased respiratory drive, reduced lung compliance, increased airway resistance, or impaired gas exchange trigger reflex activation of diaphragmatic, intercostal, and accessory muscle contraction.

\section*{Causes}
\begin{itemize}
  \item \textbf{Upper airway:} Croup, epiglottitis, foreign body aspiration, anaphylaxis, peritonsillar abscess
  \item \textbf{Lower airway:} Asthma exacerbation, bronchiolitis, COPD exacerbation, bronchiectasis
  \item \textbf{Parenchymal:} Pneumonia, ARDS, pulmonary edema (cardiogenic), interstitial lung disease, COVID-19 pneumonitis
  \item \textbf{Pleural:} Pneumothorax, pleural effusion, hemothorax
  \item \textbf{Vascular:} Pulmonary embolism, pulmonary hypertension
  \item \textbf{Other:} Metabolic acidosis (DKA, salicylate poisoning), neuromuscular disease (Guillain-Barr\'e, myasthenia gravis), chest wall trauma (flail chest), panic attack
\end{itemize}

\section*{When to See a Doctor}
See your doctor if:
\begin{itemize}
  \item Respiratory rate greater than 30 breaths/min in adults (or greater than 60 in infants), with visible distress
  \item Oxygen saturation less than 90\% room air, or rising pCO2 with acidosis
  \item Altered mental status, confusion, or somnolence (CO2 narcosis)
  \item Inability to speak full sentences (one-word dyspnea)
  \item Cyanosis, diaphoresis, or hypotensive shock
\end{itemize}

\section*{Related Symptoms}
\begin{itemize}
  \item Tachypnea, dyspnea, accessory muscle use, nasal flaring, retractions, grunting, paradoxical breathing, hypoxia
\end{itemize}
\textbf{Important:} Grunting in an infant is a compensatory mechanism that generates intrinsic PEEP to maintain alveolar patency and should not be mistaken for comfort.

\section*{Self-Care / Home Management}
\begin{itemize}
  \item Sit you upright in a tripod position (leaning forward with arms supported on a table)
  \item Pursed-lip breathing technique (4 seconds in, 6 seconds out) to reduce air trapping
  \item Activate emergency services immediately -- respiratory distress is a medical emergency
  \item If known asthma or COPD, administer rescue inhaler (albuterol 2--4 puffs) while awaiting help
  \item Loosen tight clothing and ensure adequate ventilation in the room
\end{itemize}

\section*{How Your Doctor Will Evaluate You}
\begin{itemize}
  \item Primary survey (ABCs): airway patency, breathing (rate, depth, symmetry), circulation (pulse, BP)
  \item Pulse oximetry, arterial blood gas (pH, PaO2, PaCO2, HCO3) to quantify severity
  \item Chest X-ray for pulmonary edema, pneumonia, pneumothorax, or effusion
  \item ECG and cardiac biomarkers to assess for cardiogenic causes (pulmonary edema, ischemia)
  \item CT pulmonary angiography for suspected pulmonary embolism; lung ultrasound for rapid bedside assessment
\end{itemize}

\section*{Treatment Options}
\begin{itemize}
  \item High-flow supplemental oxygen via non-rebreather mask to target SpO2 greater than 92\% (or 88--92\% in COPD)
  \item Noninvasive positive pressure ventilation (CPAP, BiPAP) for pulmonary edema or COPD exacerbation
  \item Bronchodilators (albuterol nebulized) and systemic corticosteroids for obstructive lung disease
  \item Treat underlying cause: antibiotics for pneumonia, diuretics for pulmonary edema, anticoagulation for PE
  \item Endotracheal intubation with mechanical ventilation for impending or actual respiratory failure
\end{itemize}

\clearpage
